% Section 14: Substitution, ledger transport, and strong normalization
\label{sec:normalization-body}

The proof theory of $\mathrm{TND}_0$ begins with substitution, because an ordered linear
fragment cannot borrow the usual substitution lemma: the order is part of the data, and the
ledger records uses that an ordinary substitution would silently move.  This section proves
an ordered one-hole substitution theorem with an exact ledger splice, derives conditional
cut admissibility from it, shows that raw ledger equality is the wrong target, and then
transports the whole apparatus through proof contexts to obtain strong normalization.

\subsection{Ordered one-hole substitution}

\begin{definition}[One-hole context]\label{def:one-hole-context}
A \emph{one-hole proof context} is a proof tree with a distinguished hole at a leaf.  Two
proof trees are \emph{concatenable} along the hole when the type of the conclusion at the
hole matches the type of the root of the inserted tree; concatenation is order-respecting
and is written $Q[\cdot]$.
\end{definition}

\begin{definition}[Ordered substitution]\label{def:ordered-substitution}
Let $\sigma$ be a substitution assigning to each free slot of the antecedent block a
position in a target context.  The \emph{ordered substitution} $P\sigma$ is obtained by
replacing each free occurrence in order of the context, without permutation.
\end{definition}

\begin{theorem}[Ordered substitution with ledger splice]\label{thm:substitution-splice}
Let $\Lambda_\sigma$ be the ledger change induced by the substitution $\sigma$.  For a
one-hole context $L\,Q\,R$ in which $Q$ sits between a left ledger block $L$ and a right
ledger block $R$,
\[
  \mathrm{ledger}\bigl(L\,\Lambda_\sigma\,R\bigr)
  =
  L\circ\Lambda_\sigma\circ R ,
\]
an exact splice: the left block is untouched, the middle block is transported by $\sigma$,
and the right block is untouched, with no reordering.
\end{theorem}

\begin{proof}
Induction on the context $L\,Q\,R$.  At the hole the ledger is exactly
$\Lambda_\sigma$ by Definition~\ref{def:ordered-substitution}.  For each rule above or
below the hole, the rule's own ledger contribution is a block that does not commute with
the hole's block, which is the content of linearity: the induction hypothesis is applied
to the strict subcontexts, and the blocks compose in the order fixed by the derivation.
\end{proof}

\begin{corollary}[Fresh-cut admissibility, conditionally]\label{cor:fresh-cut}
If the transported ledger $\Lambda_\sigma$ is \emph{fresh}, meaning no antecedent is
consumed twice, then the cut on the substituted formula is admissible: the composite
$L\,\Lambda_\sigma\,R$ is derivable whenever the substitutes are.
\end{corollary}

\begin{proof}
Apply Theorem~\ref{thm:substitution-splice} and check that each antecedent in the spliced
ledger occurs exactly once.  Freshness is exactly the condition that no contraction is
implicitly required, and without it the splice would duplicate an antecedent in violation
of Definition~\ref{def:tnd0}(2).
\end{proof}

\begin{remark}[The condition is not removable]
Corollary~\ref{cor:fresh-cut} is conditional and stays conditional.  Dropping freshness
would allow a ledger with a repeated antecedent, which no rule of the fragment can
introduce; the resulting "admissible" cut would have no derivation.
\end{remark}

\subsection{Raw ledger equality fails}

\begin{nogo}[Ledger equality is not the right invariant]\label{nogo:raw-ledger}
Raw equality of ledgers is not preserved by beta-like reduction: there are a proof tree and
its reduct whose ledgers are unequal as words while both record the same current authority.
Consequently the target of the proof transformation is not ledger equality but current
ledger equality plus a certificate for the retired part.
\end{nogo}

\begin{proof}
A principal-implication detour consumes the antecedent in one order and its reduct consumes
it in another; both consume it exactly once.  The ledgers therefore differ as words while
agreeing in the multiset of current antecedents.
\end{proof}

\begin{proposal}[Ledger transport, not ledger identity]\label{prop:ledger-transport}
The correct statement is a \emph{transport}: a reduction relates the pre-state ledger to the
post-state ledger by a declared transformation, and retired authority leaves the current
ledger while surviving in the certificate.  This is controlled forgetting in the sense of
Proposition~\ref{prop:three-grades}, not equality.
\end{proposal}

\subsection{Transport through contexts}

\begin{lemma}[Freshness reflects through discharge]\label{lem:freshness-reflects}
If a ledger is fresh above a discharge, then its image under the discharge is fresh, and
conversely a fresh image implies freshness above whenever the discharge consumes exactly
one antecedent.
\end{lemma}

\begin{proof}
The discharge removes exactly one antecedent occurrence and appends the discharged
assumption to the premise's right block.  Freshness is a condition on occurrences, so it is
preserved in both directions under this exact-one bookkeeping.
\end{proof}

\begin{theorem}[Root transport lifts through every context]\label{thm:root-transport}
Let
\[
  L\,[x^{b}]\,Q\,S\ \leadsto\ L\,S\,Q
\]
be the root beta-like transport of a principal detour.  Then for every one-hole proof
context $C[\cdot]$,
\[
  C\bigl(L\,[x^{b}]\,Q\,S\bigr)\ \leadsto\ C\bigl(L\,S\,Q\bigr),
\]
and the transported ledger differs from the original by the declared transformation only.
\end{theorem}

\begin{proof}
By induction on the context, using Lemma~\ref{lem:freshness-reflects} at each discharge and
Theorem~\ref{thm:substitution-splice} at each substitution.  Linearity ensures that no
context can duplicate the detour, so the transport is well defined at every depth.
\end{proof}

\begin{corollary}[Contextual closure]\label{cor:contextual-closure}
The beta-like reduction relation of $\mathrm{TND}_0$ is closed under one-hole contexts.
\end{corollary}

\begin{proof}
Immediate from Theorem~\ref{thm:root-transport}.
\end{proof}

\subsection{Strong normalization}

\begin{definition}[Proof-tree ledger and discharge count]\label{def:proof-ledger}
For a checked $\mathrm{TND}_0$ proof tree $D$ let $\Lambda_D$ be its \emph{complete ledger}:
the word of resource records of $D$ in tree order, each carrying an occurrence key, a
status, and, when discharged, the binder that discharged it.  The status is
\emph{open} for a use still owed to the ordered context and \emph{discharged} for a use
bound by $(\multimap I)$.  The open projection of $\Lambda_D$ is the ordered context of the
judgment at the root, so the complete ledger is finer than the ledger of
Definition~\ref{def:ledger-judgment}, which records only the current part.  Write
$\operatorname{dchg}(\Lambda_D)$ for the number of discharged records, and
$\operatorname{nodes}(D)$ for the number of nodes of the tree.
\end{definition}

\begin{theorem}[Exact ledger descent]\label{thm:ledger-descent}
For every certified contextual beta step $D\Longrightarrow_\beta D'$,
\begin{equation}
  |\Lambda_{D'}|=|\Lambda_D|-1,
  \qquad
  \operatorname{dchg}(\Lambda_{D'})=\operatorname{dchg}(\Lambda_D)-1 .
  \label{eq:ledger-descent}
\end{equation}
In the unfolded proof-tree representation one also has
$\operatorname{nodes}(D')=\operatorname{nodes}(D)-3$.
\end{theorem}

\begin{proof}
At the root, the transport of Theorem~\ref{thm:root-transport} replaces
\[
  \Lambda_R
  =L\circ[x^{\mathrm{discharged},b}]\circ Q\circ\Lambda_\sigma
  \quad\leadsto\quad
  \Lambda_C
  =L\circ\Lambda_\sigma\circ Q ,
\]
retiring exactly the bound target record $x$ and moving the argument block $S=\Lambda_\sigma$
across the discharged suffix $Q$; no other record is deleted, duplicated, relabelled, or
renamed, so both counts fall by exactly one.  An introduction ancestor discharges the same
surviving open record on both sides and extends both ledgers by one record at the same
position; either elimination ancestor concatenates the same sibling ledger on both sides.
Induction along the proof-position path therefore gives \eqref{eq:ledger-descent} at the
outer root.

For the node count, linearity makes substitution graft $\sigma$ exactly once: the reduct
removes the elimination node, the introduction node, and the target assumption leaf, and
nothing else.
\end{proof}

\begin{remark}[Why the measure is a ledger count and not a detour count]
Counting principal-implication detours is the tempting alternative, since a root step
contracts one.  It is not a descent measure: contracting a detour can expose another at the
same position.  In the $\lambda$-term reading of the fragment,
\[
  \bigl((\lambda f.f)(\lambda x.x)\bigr)z
  \;\longrightarrow\;
  (\lambda x.x)z
\]
has exactly one principal detour on each side; every variable is used once and no exchange
is needed.  The ledger counts of Theorem~\ref{thm:ledger-descent} are immune to this,
because they measure what a step retires rather than what it exposes.
\end{remark}

\begin{theorem}[Strong normalization]\label{thm:strong-normalization}
Every reduction sequence in $\mathrm{TND}_0$ terminates.  Indeed, a sequence starting at $D$
has length at most $\operatorname{dchg}(\Lambda_D)$.
\end{theorem}

\begin{proof}
By Theorem~\ref{thm:root-transport} a step may be applied at any depth, and by
Theorem~\ref{thm:ledger-descent} each step lowers $\operatorname{dchg}$ by exactly one.  An
infinite sequence would therefore give an infinite strictly descending sequence of natural
numbers, which is impossible.
\end{proof}

\begin{remark}[Where linearity was used]
Theorem~\ref{thm:strong-normalization} depends on linearity three times: in
Theorem~\ref{thm:substitution-splice}, which forbids reordering the ledger; in
Lemma~\ref{lem:freshness-reflects}, which keeps freshness intact through discharge; and in
the node count, where linearity makes the graft single-use.  The same argument fails for a
fragment with contraction, where one retires one record while copying a subtree.
\end{remark}

\subsection{Summary}

\begin{enumerate}
\item Ordered substitution produces an exact ledger splice
  (Theorem~\ref{thm:substitution-splice}).
\item Cut admissibility is conditional on freshness
  (Corollary~\ref{cor:fresh-cut}), and the condition is not removable.
\item Ledger equality is the wrong invariant; transport plus certificate is right
  (No-go~\ref{nogo:raw-ledger}, Proposal~\ref{prop:ledger-transport}).
\item Root transport lifts through all contexts, and reduction is strongly normalizing
  (Theorem~\ref{thm:root-transport}, Corollary~\ref{cor:contextual-closure},
  Theorem~\ref{thm:strong-normalization}).
\end{enumerate}
